% !TeX root = ../main.tex
\section{模型建立与求解}

以加权残差最小化为可编译公式示例。正式论文应给出与题目机制一致的
目标函数、约束条件、参数含义和求解方法：
\begin{equation}
  \hat{\theta}
  = \arg\min_{\theta \in \Omega}
    \sum_{i=1}^{n} w_i\left[y_i-f(x_i;\theta)\right]^2,
  \qquad w_i>0.
\end{equation}
其中，$\Omega$ 表示由物理、业务或竞赛规则确定的可行域。

正式写作前应说明基础模型、触发模型升级的证据、新增机制，以及共享参数、
条件参数和干扰参数的角色；给出可辨识性边界。只有确实承载正文证据时，才
插入自动生成的模型推理与参数角色表，不要把内部审计台账原样堆入论文。
